(1) 可逆绝热理想气体满足 \(T_2/T_1=(p_2/p_1)^{(\gamma-1)/\gamma}=\tau\)，\(T_4/T_3=1/\tau\)，故 \(T_2=\tau T_1\)、\(T_4=T_3/\tau\)。压缩与膨胀比功分别为 \(c_p(T_2-T_1)\) 与 \(c_p(T_3-T_4)\)。
(2) 理想回热器仅在 \(T_4>T_2\) 时有正回热量，即 \(T_3>\tau^2T_1\)。冷侧由 \(T_2\) 预热至 \(T_4\)，热侧由 \(T_4\) 冷却至 \(T_2\)；外部加热 \(q_{\rm in}=c_p(T_3-T_4)\)，外部排热 \(q_{\rm out}=c_p(T_2-T_1)\)。
(3) \(w_{\rm net}=c_p[T_3(1-1/\tau)-T_1(\tau-1)]\)，回热不改变压缩机与透平功。效率 \(\eta=1-q_{\rm out}/q_{\rm in}=1-\tau T_1/T_3\)；无回热时 \(\eta_0=1-1/\tau\)，在 \(T_3>\tau^2T_1\) 区间回热效率更高。
(4) \(dw_{\rm net}/d\tau=c_p(T_3/\tau^2-T_1)\)，最大值位于 \(\tau_*=\sqrt{T_3/T_1}\)，对应 \(r_{p,*}=(T_3/T_1)^{\gamma/[2(\gamma-1)]}\)。此时 \(T_4=T_2\)，回热可用温差恰为零；净功最优位于回热有效区间的边界。
